@@M0@@; let us work, for example, on @@M1@@. Then @@M2@@ is defined by the module @@M3@@, and @@M4@@ by @@M5@@, the canonical morphism sending @@M6@@ to @@M7@@ for every @@M8@@. Then @@M9@@ has as a basis the elements @@M10@@ for @@M11@@.

We have @@M12@@; put @@M13@@.
Since the sequence of the @@M14@@ is regular in @@M15@@, the same is true of the sequence @@M16@@ deduced from it by an invertible matrix, and hence also of the subsequence @@M17@@ (we are in a noetherian case). Consequently, @@M18@@ is a regular morphism.

We must therefore show that the @@M19@@ generate the Ideal of @@M20@@. In @@M21@@, for @@M22@@, we shall denote by @@M23@@ the element @@M24@@ of @@M25@@, considered as having degree @@M26@@. Over @@M27@@, @@M28@@, and @@M29@@ is defined by the morphism sending @@M30@@ to @@M31@@. Under this morphism, the image of @@M32@@ is @@M33@@. Thus the @@M34@@ belong to the Ideal of @@M35@@.

Since every element of @@M36@@ can be expressed as a polynomial in @@M37@@, to show that the @@M38@@ generate the Ideal of @@M39@@ it suffices to show that no nonzero element of @@M40@@ belongs to the Ideal of @@M41@@, which amounts to showing that the homomorphism
@@M42@@
is injective, since @@M43@@ is not a zero divisor in @@M44@@. Grading @@M45@@ by assigning degree @@M46@@ to @@M47@@ and degree @@M48@@ to the @@M49@@, the kernel of the homomorphism under consideration is a graded ideal. Now it is clear that
